\section{Discussion}\label{sec:discussion}

The central fact this paper establishes is one of scale and design, not
of behaviour: a designed, rules-based insurance system for
trade-displaced manufacturing workers would carry a gross annual cost of
the same order as the improvised support the UK already commits when a
single large employer fails---while covering every eligible displaced
worker rather than those behind the season's headline intervention. The
comparison does not say the steel interventions were wrong; they bought
industrial capacity, supply-chain security and time, not only worker
support. It says the worker-support component of UK trade policy is
priced ad hoc, delivered by postcode, and never evaluated---and that the
alternative is affordable enough to be designed on its merits.

The costings carry their limits on their face. They are gross envelopes:
the microsimulation stage will show the net Exchequer cost of each arm to
be materially lower once income tax and the Universal Credit taper claw
back payments, and will locate the incidence across the income
distribution. They hold behaviour fixed: the US evidence suggests the
wage-insurance arm in particular would change re-employment behaviour in
the direction that reduces its net cost, and the bracketed sensitivity is
no substitute for UK evidence that does not yet exist. Eligibility is
bounded rather than identified: any real scheme must choose between a
certification mechanism it would have to build and a breadth it would
have to fund. And the flows are cyclical: an insurance scheme priced on
recent redundancy rates must expect to pay out multiples of its
steady-state cost in the downturn that is precisely when it is needed.

Three implications follow for the policy conversation this paper is
meant to start. First, the cheapest arm on the table is a rule change,
not a programme: disregarding statutory redundancy pay from the
Universal Credit capital test repairs an interaction in which the
compensation Parliament mandates for job loss switches off the safety
net Parliament provides for it---an incoherence, not a policy choice,
and one whose cost the microsimulation stage can pin down. Second, the
government's own proposed Unemployment Insurance is most of an
adjustment instrument already; the marginal questions are its duration
and whether trade-displaced workers warrant the extended horizon the US
has offered for sixty years. Third, wage insurance is the arm the UK
debate has never priced in its modern form, the one with the strongest
causal evidence behind it, and the one whose UK-specific interaction with
means-testing makes the imported evidence least sufficient---the natural
centrepiece for the full microsimulation stage. Two of this paper's
results sharpen its case: the statutory break-even is
\TaaBreakEvenMonths\ months of faster re-employment per year of
entitlement (\TaaBreakEvenMonthsFull\ over the full two years, at the
top of the range the US evidence observes); and the UK's own historical
design---ISERBS extended from steel to all manufacturing---would cost
\pounds\TaaIserbsCost\,million a year, the same cost class as the
US-style arm, giving a precedent-anchored starting point that requires
no imported blueprint.

The paper's discipline throughout has been to separate what is known
(statutory parameters, redundancy flows, earnings distributions), what
is designed (the four arms), what is bracketed (take-up, re-employment,
penalties), and what awaits the microdata (net costs, incidence, the
capital-gate repair). Within those separations, the finding stands: the
United Kingdom does not lack the money for trade adjustment assistance.
It has been spending it, one emergency at a time, without the
instrument.
